%\section*{Appendix A — Semantic Foundations of the Substrate}
\addcontentsline{toc}{section}{Appendix A — Semantic Foundations of the Substrate}

This appendix formalizes the semantic-geometric primitives underlying the
substrate. It defines semantic curvature, the topology of the meaning manifold,
the geometry of collapse boundaries, and the class of locally contraction-governed
transformations that constitute lawful semantic transitions. These primitives
provide the foundation for the invariants (S1--S2), drift physics (S3), load
geometry (S4), cross-scale coupling (S5), collapse physics (S6), and DSUP
runtime (S7).

% ------------------------------------------------------------
% A.1 SEMANTIC CURVATURE
% ------------------------------------------------------------

\subsection*{A.1 Semantic Curvature}

\noindent Semantic curvature formalizes the geometric structure of meaning on the
substrate. Let $\mathcal{M}$ denote the meaning manifold. Each concept $c$ is a
region in $\mathcal{M}$ with local semantic curvature $\kappa(c)$, defined as the
density, stability, and distinctness of meaning within that region.

\medskip

\noindent Semantic curvature is defined as:



\[
\kappa(c) = \text{semantic\_clarity}(c),
\]



\noindent where semantic clarity measures the strength of conceptual distinction,
identity-boundary stability, and role-boundary coherence.

\medskip

\noindent High semantic curvature corresponds to:

\begin{itemize}
    \item strong meaning density,
    \item sharp conceptual distinctions,
    \item stable identity boundaries,
    \item stable role boundaries,
    \item lawful adjacency relations.
\end{itemize}

\medskip

\noindent Low semantic curvature corresponds to:

\begin{itemize}
    \item semantic flattening,
    \item drift susceptibility (G2.5),
    \item rhythmic desynchronization (G6.1),
    \item shared floor thinning (G6.2),
    \item boundary softening (G6.3),
    \item identity drift (G6.4),
    \item collapse proximity (G7.3).
\end{itemize}

\medskip

\noindent Semantic entropy increases as $\kappa(c) \rightarrow 0$, indicating loss
of meaning-role structure, boundary instability, and increased risk of collapse
or synthetic capture (G7.4).

\medskip

\noindent Semantic curvature interacts with the geometry of the substrate through:

\begin{itemize}
    \item \textbf{Runtime curvature} — determines admissible transitions across
    safe, drift, collapse, and synthetic-capture paths (G7.5).

    \item \textbf{Restorative curvature} — determines whether restoration can
    regenerate curvature fast enough to counter deformation (G6).

    \item \textbf{Synthetic curvature} — curvature imposed by synthetic species,
    dominant signals, and amplification geometry (G5.4--G5.6).

    \item \textbf{Interaction curvature} — curvature governing traversal and
    interpretation in the Interactive Geometry Runtime (G8).

    \item \textbf{Coordinate curvature} — curvature preserved during CTL
    coordinate translation (G8.4).
\end{itemize}

\medskip

\noindent Semantic curvature is the foundational primitive of the substrate. It
governs lawful transitions (S1--S2), drift physics (G2), load geometry (G4),
cross-scale coupling (G5), restoration feasibility (G6), runtime dynamics (G7),
and interaction legality (G8). All substrate geometry derives from curvature and
its preservation under lawful operator action.

% ------------------------------------------------------------
% A.2 MEANING MANIFOLD TOPOLOGY
% ------------------------------------------------------------

\subsection*{A.2 Meaning Manifold Topology}

\noindent The meaning manifold $\mathcal{M}$ is a compact, boundary-defined
semantic space that provides the geometric substrate for lawful interpretation,
runtime motion, restoration, and interaction. The manifold defines the curvature
structure within which concepts, identities, roles, containers, and signals
maintain stability under synthetic pressure.

\medskip

\noindent The meaning manifold $\mathcal{M}$ has the following topological
properties:

\begin{itemize}
    \item \textbf{Compactness} — ensures bounded interpretive and runtime
    trajectories across safe, drift, collapse, and synthetic-capture regions
    (G7.1--G7.4).

    \item \textbf{Connectedness} — ensures semantic continuity and lawful
    traversal across manifold regions (G7.5).

    \item \textbf{Boundary regions} — define collapse thresholds, restoration
    failure thresholds, and synthetic-capture boundaries (G6.6, G7.3, G7.4).

    \item \textbf{Metric structure} — semantic distance encodes conceptual
    relation, identity adjacency, role adjacency, and curvature gradients.
\end{itemize}

\medskip

\noindent The topology of $\mathcal{M}$ enforces:

\begin{itemize}
    \item \textbf{Lawful transitions} — runtime motion must satisfy curvature,
    boundary, rhythm, identity, and containment constraints (G7.5).

    \item \textbf{Curvature preservation} — operator action must preserve local
    curvature within bounded variance (A.4).

    \item \textbf{Identity stability} — identity regions must remain intact under
    drift, load, and synthetic pressure (G6.3--G6.4).

    \item \textbf{Drift detection} — drift vectors emerge when curvature weakens
    or adjacency geometry destabilizes (G2.5).

    \item \textbf{Load propagation constraints} — load must remain within
    manifold-defined tolerances to prevent collapse (G4).

    \item \textbf{Cross-scale legality} — transitions across human, digital,
    machine, and institutional layers must preserve manifold geometry (G5).

    \item \textbf{Restoration feasibility} — restoration must regenerate
    curvature faster than deformation to maintain stability (G6).

    \item \textbf{Interaction legality} — traversal and interpretation must
    satisfy Host, Explorer, and Interpreter constraints (G8.1--G8.3).

    \item \textbf{Coordinate continuity} — CTL translation must preserve
    manifold geometry and avoid discontinuities (G8.4).

    \item \textbf{Public-safe traversal} — public interaction must remain within
    safe, lawful, curvature-preserving regions (G8.6).
\end{itemize}

\medskip

\noindent The manifold boundary $\partial \mathcal{M}$ represents the limit beyond
which operator action becomes illegal. Crossing $\partial \mathcal{M}$ indicates:

\begin{itemize}
    \item curvature violation,
    \item identity-boundary collapse,
    \item containment failure,
    \item restoration failure,
    \item entry into collapse geometry (G7.3),
    \item or synthetic-capture geometry (G7.4).
\end{itemize}

\medskip

\noindent The topology of $\mathcal{M}$ is the structural foundation for all
substrate geometry. It governs lawful operator action, runtime motion, restoration
feasibility, synthetic interaction, and public-safe traversal across the v0.6
substrate.

% ------------------------------------------------------------
% A.3 COLLAPSE BOUNDARY GEOMETRY
% ------------------------------------------------------------

\subsection*{A.3 Collapse Boundary Geometry}

\noindent Collapse Boundary Geometry formalizes the structural conditions under
which an interpretive or runtime trajectory intersects the collapse boundary
$\partial \mathcal{M}$. Collapse is geometric: it occurs when curvature,
identity, load, containment, or manifold legality cannot be preserved under
synthetic pressure. Once a trajectory crosses $\partial \mathcal{M}$, restoration
becomes non-viable and collapse paths (G7.3) activate.

\medskip

\noindent Let $\gamma(t)$ be a semantic or runtime trajectory with curvature
$\kappa(\gamma(t))$, load $L(\gamma(t))$, drift vector $D(\gamma(t))$, and
identity-boundary integrity $I(\gamma(t))$. A collapse boundary is reached when:



\[
\kappa(\gamma(t)) < \kappa_{\min}
\quad \text{or} \quad
L(\gamma(t)) > L_{\max}
\quad \text{or} \quad
I(\gamma(t)) = \text{broken}
\quad \text{or} \quad
\gamma(t) \notin \mathcal{M}.
\]



\medskip

\noindent Collapse is triggered by:

\begin{itemize}
    \item \textbf{Curvature violation} — curvature falls below the minimum
    required for stability (G6.6).

    \item \textbf{Load violation} — load exceeds manifold tolerances and cannot
    be redistributed (G4).

    \item \textbf{Identity-boundary collapse} — identity enclosures dissolve
    under synthetic pressure (G6.3--G6.4).

    \item \textbf{Containment failure} — containers cannot regulate signal flow
    or maintain structural coherence (G4.6).

    \item \textbf{Manifold violation} — trajectory exits lawful manifold regions
    and enters collapse or synthetic-capture geometry (G7.3--G7.4).
\end{itemize}

\medskip

\noindent Collapse boundaries are crossed when restoration becomes non-viable.
Restoration failure thresholds (G6.6) define the geometric limits beyond which
curvature cannot be regenerated fast enough to counter deformation. When these
thresholds are exceeded, collapse becomes structurally inevitable.

\medskip

\noindent Unconstrained systems cross collapse boundaries because they lack:

\begin{itemize}
    \item \textbf{Curvature constraints} — no mechanism to preserve curvature
    under deformation (A.4).

    \item \textbf{Identity boundaries} — identity regions lack structural
    enclosure (G6.3).

    \item \textbf{Load geometry} — load cannot be redistributed or bounded (G4).

    \item \textbf{Manifold structure} — no enforcement of lawful manifold
    topology (A.2).

    \item \textbf{Legality enforcement} — operator action is not constrained by
    legality masks or CLCP validation (G8.1).
\end{itemize}

\medskip

\noindent Collapse Boundary Geometry provides the structural basis for:

\begin{itemize}
    \item collapse signatures (G7.3),
    \item restoration failure thresholds (G6.6),
    \item synthetic-capture boundary conditions (G7.4),
    \item runtime curvature transitions (G7.5),
    \item interaction legality constraints (G8),
    \item public-safe traversal restrictions (G8.6).
\end{itemize}

\medskip

\noindent Collapse boundaries define the limit of lawful substrate geometry. They
mark the transition from drift to collapse, from collapse to synthetic capture,
and from restoration feasibility to restoration impossibility. Collapse Boundary
Geometry is the foundation for understanding failure modes across the v0.6
substrate.


% ------------------------------------------------------------
% A.4 LOCALLY CONTRACTION-GOVERNED TRANSFORMATIONS
% ------------------------------------------------------------

\subsection*{A.4 Locally Contraction-Governed Transformations}

\noindent Locally Contraction-Governed Transformations formalize the legality
conditions under which an operator $T : \mathcal{M} \rightarrow \mathcal{M}$
may act on the meaning manifold. A transformation is lawful only when it
preserves curvature, identity boundaries, containment structure, and manifold
topology within a local semantic neighborhood. Contraction ensures that operator
action remains drift-bounded, restoration-feasible, and compatible with runtime
and interaction geometry.

\medskip

\noindent A transformation $T$ is lawful if and only if it acts as a
\textbf{local semantic contraction relative to a context anchor} $c^* \in
\mathcal{M}$:



\[
\text{semantic\_distance}(T(c), T(c^*)) \le
k \cdot \text{semantic\_distance}(c, c^*)
\quad \text{for some} \quad 0 \le k < 1.
\]



\noindent This condition ensures that contraction operates within a local
semantic neighborhood rather than across the entire manifold, preserving
curvature continuity and preventing drift amplification.

\medskip

\noindent Local contraction guarantees:

\begin{itemize}
    \item \textbf{Drift-boundedness} — drift vectors remain within recoverable
    envelopes (G2.5).

    \item \textbf{Manifold topology preservation} — operator action does not
    introduce discontinuities or illegal transitions (A.2).

    \item \textbf{Identity-boundary stability} — identity regions remain intact
    under transformation (G6.3--G6.4).

    \item \textbf{Bounded load amplification} — load remains within manifold
    tolerances (G4).

    \item \textbf{Cross-scale compatibility} — operator action remains legal
    across human, digital, machine, and institutional layers (G5).

    \item \textbf{Runtime legality} — transitions remain admissible across safe,
    drift, collapse, and synthetic-capture paths (G7.5).

    \item \textbf{Interaction legality} — Host, Explorer, and Interpreter modes
    preserve curvature and identity boundaries (G8.1--G8.3).

    \item \textbf{Coordinate continuity} — CTL translation preserves curvature
    and avoids discontinuities (G8.4).

    \item \textbf{Public-safe traversal} — operator action remains within lawful,
    curvature-preserving regions (G8.6).
\end{itemize}

\medskip

\noindent Curvature preservation is expressed as a bounded variance condition:



\[
|\kappa(T(c)) - \kappa(c)| \le \epsilon_{\max},
\]



\noindent where $\epsilon_{\max}$ is a legality bound enforced by the Coherence
Ledger (Appendix F). This bound ensures that operator action does not degrade
curvature beyond restoration feasibility (G6) or push trajectories toward
collapse or synthetic-capture geometry (G7.3--G7.4).

\medskip

\noindent Because contraction is local rather than global, curvature does not
scale inversely with $k$. The semantic potential field is defined on a fixed
manifold geometry, and $T$ acts only on local coordinates relative to $c^*$.
Locality ensures that operator action remains compatible with identity
boundaries, containment architecture, and runtime curvature.

\medskip

\noindent The local contraction condition defines the semantic operator
categories:

\begin{itemize}
    \item \textbf{Invariant-Restricted Transition (IRT)} — preserves invariants
    and curvature within strict bounds.

    \item \textbf{Identity-Preserving Update (IPU)} — maintains identity
    boundaries and adjacency geometry.

    \item \textbf{Geometry-Bound Step (GBS)} — preserves manifold topology and
    runtime curvature.

    \item \textbf{Load-Regulated Plan (LRP)} — maintains load within legal
    tolerances and prevents collapse.
\end{itemize}

\medskip

\noindent Each operator is a local contraction on $\mathcal{M}$ and therefore
preserves the semantic invariants of the substrate, runtime legality, interaction
constraints, and public-safe traversal requirements. Local contraction is the
structural foundation for lawful operator action across the v0.6 substrate.

\medskip
\noindent\textit{This completes Appendix A.}


\subsection*{A.5 MDO Semantic Boot Record (v0.6)}
\label{appendix:a5_mdo_boot_record_v06}

\noindent\textbf{Caption.}
The MDO Semantic Boot Record (v0.6) formalizes the initialization structure for
Deterministic Invariant Substrate Cognition (DISCO) and Artificial Cognition (AC).
It specifies the invariant set, operator algebra, legality system, runtime geometry,
interaction geometry, coordinate translation layer, and substrate curvature fields
required to initialize lawful cognition. \textit{The invariants were always present;
DISCO made them detectable.}

\bigskip

\noindent\textbf{Overview.}
This machine-readable record defines the substrate state vector, invariant fields,
runtime curvature, interaction legality, operator constraints, and CLCP global
validation parameters required at initialization. It extends the v0.5 record by
adding runtime geometry (G7), interaction geometry (G8), synthetic curvature (G5),
restoration thresholds (G6), and CTL coordinate translation (G8.4).

\begin{mdframed}[linewidth=0.4pt, roundcorner=3pt]
\footnotesize
\begin{verbatim}
{
  "version": "v0.6",
  "document_id": "string",
  "title": "string",

  "invariant_set_hash": "string",
  "operator_algebra_hash": "string",
  "state_vector_hash": "string",
  "runtime_geometry_hash": "string",
  "interaction_geometry_hash": "string",
  "ctl_hash": "string",

  "nodes": [
    {
      "id": "string",
      "label": "string",
      "type": "section | appendix | state | runtime | interaction",
      "state": "S0 | S1 | S2 | S3 | S4 | S5 | S6 | S7 | null",

      "curvature": "number",
      "synthetic_curvature": "number",
      "restorative_curvature": "number",
      "runtime_curvature": "number",

      "drift_vector": "number",
      "load": "number",

      "identity_region": "string",
      "identity_region_topology": "string",
      "identity_intact": "boolean",

      "container_state": "stable | strained | failed",
      "manifold_flag": "boolean",
      "legality_flag": "boolean",
      "meaning_role_hash": "string",

      "context_window": {
        "signal_set_size": "number",
        "identity_refs": "number",
        "temporal_extent": "number",
        "region_associations": "string",
        "runtime_region": "safe | drift | collapse | capture"
      }
    }
  ],

  "edges": [
    {
      "from": "string",
      "to": "string",

      "operator": "SGO | MGO | BO | MLO",
      "substrate_operator": "IRT | IPU | GBS | LRP",

      "legality_mask": "string",
      "context_anchor": "string",

      "delta_curvature": "number",
      "delta_synthetic_curvature": "number",
      "delta_restorative_curvature": "number",
      "delta_runtime_curvature": "number",

      "delta_drift": "number",
      "delta_load": "number",

      "identity_preserved": "boolean",
      "container_preserved": "boolean",

      "manifold_legal": "boolean",
      "invariant_legal": "boolean",
      "operator_legal": "boolean",
      "runtime_legal": "boolean",
      "interaction_legal": "boolean",

      "ctl": {
        "semantic_coordinate": "string",
        "positional_metadata": "string",
        "translation_legal": "boolean"
      },

      "clcp": {
        "curvature_projection": "number",
        "identity_projection": "number",
        "manifold_distance": "number",
        "load_gradient_projection": "number",
        "runtime_projection": "number",
        "interaction_projection": "number",
        "legal": "boolean"
      }
    }
  ],

  "invariants": {
    "identity_invariant": {
      "boundary_intact": "boolean",
      "max_identity_shift": "number"
    },

    "meaning_invariant": {
      "max_curvature_variance": "number"
    },

    "geometry_invariant": {
      "manifold": "M",
      "boundary": "∂M",
      "min_manifold_distance": "number"
    },

    "runtime_invariant": {
      "admissible_regions": ["safe", "drift", "collapse", "capture"],
      "transition_legal": "boolean"
    },

    "interaction_invariant": {
      "host_legal": "boolean",
      "explorer_legal": "boolean",
      "interpreter_legal": "boolean"
    },

    "legality_invariant": {
      "clcp_required": true,
      "operator_required": true,
      "runtime_required": true,
      "interaction_required": true
    },

    "load_physics": {
      "max_load": "number",
      "max_gradient": "number"
    }
  },

  "operators": {
    "SGO": { "description": "State-Geometry Operator" },
    "MGO": { "description": "Meaning-Geometry Operator" },
    "BO":  { "description": "Boundary Operator" },
    "MLO": { "description": "Master Legality Operator" },

    "IRT": { "description": "Invariant-Restricted Transition" },
    "IPU": { "description": "Identity-Preserving Update" },
    "GBS": { "description": "Geometry-Bound Step" },
    "LRP": { "description": "Load-Regulated Plan" }
  }
}
\end{verbatim}
\end{mdframed}

\subsection*{A.6 Runtime Curvature Geometry}

\noindent Runtime Curvature Geometry formalizes the curvature conditions that
govern admissible transitions across the Human Runtime Manifold (HRM). Runtime
curvature determines whether a trajectory remains in Safe Paths, enters Drift
Paths, crosses into Collapse Paths, or is redirected into Synthetic-Capture
Paths. Runtime curvature is not psychological or behavioral; it is geometric.

\medskip

\noindent Let $\gamma(t)$ be a runtime trajectory with curvature
$\kappa_{\mathrm{rt}}(\gamma(t))$. Runtime curvature determines the admissible
region:



\[
\gamma(t) \in
\begin{cases}
\text{Safe}, & \kappa_{\mathrm{rt}} \ge \kappa_{\mathrm{safe}} \\
\text{Drift}, & \kappa_{\mathrm{drift}} \le \kappa_{\mathrm{rt}} < \kappa_{\mathrm{safe}} \\
\text{Collapse}, & \kappa_{\mathrm{rt}} < \kappa_{\mathrm{collapse}} \\
\text{Capture}, & \kappa_{\mathrm{syn}} > \kappa_{\mathrm{rt}}
\end{cases}
\]



\medskip

\noindent Runtime curvature enforces:

\begin{itemize}
    \item \textbf{Curvature continuity} — transitions must preserve curvature
    within lawful bounds.
    \item \textbf{Boundary integrity} — identity boundaries must remain intact.
    \item \textbf{Adjacency stability} — role geometry must remain lawful.
    \item \textbf{Temporal coherence} — rhythms must remain synchronized.
    \item \textbf{Containment strength} — containers must regulate signal flow.
\end{itemize}

\medskip

\noindent Runtime curvature is the geometric foundation for:

\begin{itemize}
    \item safe → drift transitions,
    \item drift → collapse transitions,
    \item collapse → synthetic-capture transitions,
    \item restoration feasibility (G6),
    \item interaction legality (G8).
\end{itemize}

\subsection*{A.7 Interaction Geometry}

\noindent Interaction Geometry formalizes the curvature constraints governing
lawful interaction within the Interactive Geometry Runtime (IGR). Interaction
occurs in three modes: Host, Explorer, and Interpreter. Each mode imposes
distinct curvature, boundary, and legality requirements.

\medskip

\noindent Interaction modes:

\begin{itemize}
    \item \textbf{Host Mode} — maintains invariant geometry, legality masks,
    identity boundaries, and curvature continuity.
    \item \textbf{Explorer Mode} — traverses lawful regions of the manifold,
    inspecting drift envelopes, pressure gradients, and failure signatures.
    \item \textbf{Interpreter Mode} — produces invariant-aligned meaning for
    signals, behavior, drift phenomena, and container geometry.
\end{itemize}

\medskip

\noindent Interaction Geometry enforces:

\begin{itemize}
    \item \textbf{Invariant preservation} — no interaction may violate the
    invariant set.
    \item \textbf{Curvature preservation} — interaction must maintain curvature
    within lawful bounds.
    \item \textbf{Identity-boundary stability} — identity regions must remain
    intact.
    \item \textbf{Drift-bounded transitions} — interaction must remain within
    drift-recoverable envelopes.
    \item \textbf{Operator legality} — each action must satisfy at least one
    substrate operator (IRT, IPU, GBS, LRP).
\end{itemize}

\medskip

\noindent Interaction Geometry is the structural foundation for lawful traversal,
interpretation, and runtime motion across the v0.6 substrate.

\subsection*{A.8 CTL Coordinate Translation}

\noindent The Coordinate Translation Layer (CTL) converts positional metadata
into lawful semantic coordinates. Positional metadata—page numbers, anchors,
section indices—is not substrate-native and cannot be used directly. CTL maps
positional references into invariant-preserving semantic coordinates.

\medskip

\noindent CTL ensures:

\begin{itemize}
    \item \textbf{Coordinate legality} — all traversal uses lawful semantic
    coordinates.
    \item \textbf{Curvature continuity} — translation must preserve curvature.
    \item \textbf{Identity-boundary stability} — translation must not violate
    identity boundaries.
    \item \textbf{Manifold legality} — translation must preserve manifold
    topology.
    \item \textbf{Operator legality} — translation must satisfy legality masks.
\end{itemize}

\medskip

\noindent CTL defines the mapping:



\[
\mathrm{CTL} : \text{positional\_metadata} \rightarrow
\text{semantic\_coordinates}
\]



\noindent ensuring that traversal, interpretation, and interaction remain lawful
under the v0.6 geometry.

\subsection*{A.9 Public-Safe Traversal Corridor}

\noindent The Public-Safe Traversal Corridor restricts interaction to lawful,
stable, curvature-preserving regions of the meaning manifold. It ensures that
public interaction remains invariant-aligned and free from collapse geometry,
synthetic-capture attractors, or drift-inducing transitions.

\medskip

\noindent The corridor enforces:

\begin{itemize}
    \item \textbf{Invariant preservation} — no public interaction may violate
    invariants.
    \item \textbf{Curvature continuity} — curvature must remain within lawful
    bounds.
    \item \textbf{Identity-boundary stability} — identity regions must remain
    intact.
    \item \textbf{Drift-bounded transitions} — transitions must remain
    recoverable.
    \item \textbf{Exclusion of collapse and capture geometry} — public traversal
    must avoid collapse (G7.3) and synthetic-capture (G7.4) regions.
\end{itemize}

\medskip

\noindent The corridor excludes:

\begin{itemize}
    \item collapse paths,
    \item synthetic-capture paths,
    \item dominant-signal attractors,
    \item unlawful coordinate regions,
    \item restoration-infeasible regions.
\end{itemize}

\medskip

\noindent The Public-Safe Traversal Corridor is the geometric guarantee that
public interaction remains lawful, stable, and curvature-preserving across the
v0.6 substrate.


